% 【研读实验室】所属：第9章（由 chapters/revised/09-comparative.tex \input 引入）
\minitopic{研读实验室：比较是一种受约束的重构}\index{内部解释}\index{结构比较}\index{规范建构}\index{不可通约} 比较认识论不是把两份术语表并排。研究者先在每个文本内部说明问题、论证和实践角色，再建立可检验的比较维度。若一开始便规定所有传统都必须回答“知识的必要充分条件是什么”，会把不以定义为中心的材料裁成缺失版本。 受约束重构同时避免另一个极端：因历史语境不同便拒绝任何跨传统判断。概念无需完全同一才能比较。只要说明比较轴，例如错误如何被发现、权威如何可纠正、主体怎样形成能力，就能分析结构相似与差异。不可直接翻译本身也是结果。 比较结论应有强度层级。“两段文本都讨论学习”最弱；“它们对学习中反思与实践的关系给出结构相似答案”较强；“它们表达同一理论”最强，需要语义、推论与历史证据。写作时明确强度，防止从主题相近滑到理论相同。

\minitopic{文本证据的四层}第一层是版本与篇章。古典作品可能有编纂层、异文和注释传统；引用必须给出篇章或段落，而不只作者名。第二层是词语在句中的语法和语义，不能以一个现代译词固定全部用法。第三层是论证角色：该句是提问者立场、寓言声音、反面例子还是作者认可。第四层是实践背景：概念用于辩论、修养、治理或解脱，会改变理论意义。 二手研究帮助识别争议，却不能让研究者跳过原文。若两个译本差异影响论证，应并列说明并给出选择理由。读者不懂原语时，可以诚实依赖学术译本，同时降低对细微词义的断言强度。 引用名句尤其危险。以“知之为知之”为例，脱离上下文后，它很容易被包装成任何现代观点的格言\parencite[2.17]{analects2003}。“吾生也有涯”也必须放回原篇的论述进程\parencite[chap.~3]{zhuangzi2009}。先复述前后问题，再说明名句承担的推理功能；若只提供道德启发，应标为启发性使用，不伪装成历史解释。

\minitopic{六种比较关系}同源关系要求历史传播或共同文本证据；平行关系只说在无已知影响下出现相似问题；对比关系用差异照亮各自承诺；翻译关系研究概念如何跨语言移动；批判关系让一方揭示另一方盲点；建构关系从多种资源提出新规范方案。六者需要不同证据。 若没有传播证据，就不应说某思想“影响”另一传统；年代更早也不证明后者抄袭。平行比较不需要同源，却要避免竞赛式“谁先发现”。哲学价值来自论证和问题重构，不来自时间冠军。 批判比较应双向。用当代证据主义检验古典权威时，也让古典修养论追问证据主义如何形成诚实、耐心和可教的主体。若只有一方提供标准，另一方只做材料，比较仍以单一中心组织。 建构方案必须标明作者责任。研究者从荀子学习论、德性认识论和社会纠错提出新模型，这个模型不再等同任何原传统。清楚区分“文本说”“可重构为”“我据此主张”，既保护历史准确，也允许哲学创造。

\minitopic{概念簇而非一对一词典}“知”可涉及识别、会做、理解、意识到、具有智慧或道德上真正体认；“knowledge”也在不同理论中与事实性、信念、能力和断言相连。一对一翻译会让两边内部差异消失。较好的方法是建立概念簇：列出典型对象、推论、反例、主体和实践后果。 翻译可以用测试句检验。若把某处“知”译为命题知识，能否自然接宾语？是否允许真假评价？文本是否把能言与能行区分？若译为技能，是否要求重复表现？没有单一测试决定，但多项证据能限制选择。 不可译词不应成为神秘特权。研究者仍要说明它与哪些词对立、何时适用、什么算误用。音译保留开放性，却不免除分析。反过来，目标语言已有相近词也不保证理论角色相同。 概念工程还需公开规范目标。为了跨文化课堂，可设计一个宽泛“认识成就”上位词，容纳知识、理解与工夫；这有教学价值，但不是发现所有传统暗中共享一个本质。新词的便利与历史主张要分开。

\minitopic{案例一：庄周梦蝶与梦论证}两段材料都可引出梦醒区分，却处于不同问题进程。笛卡尔逐步扩大怀疑，寻找不可怀疑起点；庄周梦蝶嵌入物化、视角与身份转换的讨论\parencite{descartes1996,zhuangzi2009}。把二者都称“外部世界怀疑论”会损失终点差异。 可比较轴包括：经验内部是否有决定性梦醒标志；主体身份如何保持；怀疑的治疗目标是什么。笛卡尔路线可能把不确定性当作需克服的威胁，庄子文本则可能改变我们对固定分判的执着。这里用“可能”是因为庄子解释有争议，不能用一个寓言消灭全部诠释。 反向压力也有价值。庄子可追问寻求绝对基础是否由特定认识欲望推动；笛卡尔可追问视角转换是否仍需某些稳定判断才能表达。比较不是选胜者，而是让两种怀疑实践的隐含目标显现。

\minitopic{案例二：荀子解蔽与认识偏差}“解蔽”可与确认偏误、选择性注意和群体极化对话，因为都关心主体如何被一端遮蔽\parencite{xunzi2014}。但现代偏差研究通常以统计表现或规范模型为基准，荀子讨论还涉及欲望、名辨、学习和政治秩序。二者的错误单位并不相同。 当代理论可细分机制：信息采样偏、更新权重偏、记忆偏或表达激励偏。荀子式视角则强调长期学习、师法和心的统摄，提醒纠错不是一次提供更多事实。结合可提出多层干预：改变信息环境、训练辨别、设置批评关系并校正欲望。 建构结论应归于研究者：“我们可据此提出一种生态纠偏模型”，而非“荀子已经发现现代认知偏差”。前者承认资源来源和新组合，后者会把历史文本变成现代科学预告。

\minitopic{案例三：知行与厚知识}若“知”被规定为任何真命题的知识，“知道三角形内角性质”显然不要求以道德行动表现。某些儒学语境中的知却把体认、倾向和实践完成纳入评价\parencite{tiwald2014}。分歧可能不是事实争议，而是概念厚度不同。 比较可区分最低命题知识、能力知识、反思知识和德性化体认。厚概念能解释为何口头赞成道德原则却持续违背被批评为“不真知”；代价是难以保留日常所说“明知故犯”。薄概念能描述意志失败，却可能低估学习对人格的要求。 不必选一个词义消灭另一方。理论可承认层级：主体在薄意义知道规范，在厚意义尚未形成稳定认识成就。关键是避免在论证中从薄义前提滑到厚义结论。

\minitopic{案例四：佛教量论与当代证据}印度佛教量论传统对知识来源、知觉、推理与错误有高度技术化讨论；不同作者与后继解释差异显著\parencite{dunne2004}。把“量”直译成证据或知识来源只是起点，不能据此宣布它等同可靠主义、基础主义或贝叶斯主义。 比较需要具体文本和论证：认知是否必须新颖且不欺诳，知觉如何排除概念构造，推理标志如何成立，修行目标怎样影响认识分类。若教材只有二手概述，应降低历史断言，并把进一步原典研究列为开放任务。 当代理论可借此反思默认分类是否过度围绕命题信念；佛教材料也可受当代反事实与社会知识问题追问。对话的边界同样重要：解脱目标、心识理论和论辩语境不能从几个共享术语中删去。

\minitopic{案例五：认识自主与关系}自由主义图景常把自主理解为个人能反思并选择信念；关系性观点强调，语言、能力和自信都在依赖中形成。儒家学习关系可提示受教、劝谏和角色责任的认识意义；当代社会认识论则揭示关系也会压制、沉默和分配不公。 关系本身既非自主的敌人，也非保证。判断要看主体能否获得替代来源、提出异议、理解权威范围并在受损时退出；也要看照护和教育是否提供发展能力的真实条件。孤立可能减少支配，却也切断知识资源。 比较结论可以是规范性的：认识自主是管理不可避免依赖的能力与制度条件。这个结论受到多种传统启发，却不应倒签给任何单一古典作者。

\minitopic{研究者位置与中心主义}研究者选择问题、译本和比较轴，不可能完全无位置。位置说明不是以身份代替论证，而是公开材料能力、语言限制和规范兴趣。例如，只能阅读译本就应说明；以当代知识分析为起点，就要主动检查它是否压缩其他问题。 去西方中心化不等于把“西方”整体拒绝，也不等于对非西方材料取消批评。真正对称要求双方都有内部差异、可犯错性与论证责任。若一方被描写成抽象理性，另一方被浪漫化为整体智慧，仍是失真二分。 课程编排也表达中心。若比较只在末章出现，读者可能把前十三章当普遍理论、最后一章当文化补充。本书在各章加入比较问题，末章任务是反思方法与汇总差异。贯穿并不要求每章平均分配传统，而要求主要问题不被单一谱系垄断。

\minitopic{结课综合：建立自己的认识论地图}完成全书后，不应只选择“我是可靠主义者”。先列出不同评价：知识、确证、理解、责任、群体表现和行动许可；再说明每项由哪些条件决定。单一理论可解释部分维度，多元理论必须交代维度间关系。 其次选择三个压力案例：一个Gettier案例、一个社会或数字系统案例、一个跨传统概念案例。检验自己的立场是否在三个领域使用一致原则，若不一致，说明差异来自对象还是临时修补。理论成熟表现为知道代价，而非声称没有反例。 最后写出纠错程序：什么证据会改变立场，遇到专家分歧如何行动，如何记录不确定性，怎样避免身份偏见和来源重复。认识论不是获得一组永远正确答案，而是构造能在错误中继续学习的规范。

\begin{tcolorbox}[breakable,colback=white,colframe=blue!30!black,title={结课研读任务},fonttitle=\bfseries]
选择两个具体文本和一个当代认识论问题，分别完成内部解释、结构比较和规范建构；为每项主张标注文本证据强度；列出一个相似点、两个不可约差异和一项双向批评；最后说明新立场属于谁，以及何种反例会迫使你修正。
\end{tcolorbox}
